\section{A Distinct Almost-Prime Program}
\label{sec:almost-prime-program}

Requiring both endpoints of a square-safe 2-gap to be prime reaches the
twin-prime boundary. A separate program relaxes the second endpoint to have at
most two prime factors. That program has different local factors and a
different Type-I/Type-II formulation; it is developed separately.

Its success would not prove a surviving 2-gap or infinitely many twin primes.

\subsection{Why This Is Hard: The Type-II Barrier}
\label{sec:type-two-barrier}

The exhaustion argument and the fixed-seed scale conflict together explain
why this project's complete-period algebra cannot finish the twin-prime
program. The deeper reason the program is hard at all comes from the
contemporary theory of prime-producing sieves.

Ford and Maynard study nonnegative sequences whose difference from a
comparison model satisfies Type-I and Type-II estimates. Type-I controls
divisibility averages over many factor scales; Type-II controls bilinear
sums against arbitrary bounded coefficient sequences. Their framework proves
that a \emph{substantial Type-II range is genuinely necessary} to guarantee a
nontrivial prime lower bound: very strong Type-I information alone can still
be consistent with a sequence containing no primes.

This project's exact residue and CRT formulas are algebraic input for a
possible Type-I analysis. They are not yet a Type-I theorem over the
required short intervals, because the accumulated discrepancy norm has not
been bounded---the fixed-seed scale conflict of Section~10.1 is one face of this.
No arbitrary-coefficient Type-II estimate is currently proved for the
sieve-sequence weights.

The relaxed almost-prime program of Section~11 makes the Type-II barrier concrete.
Its scalar-centered final weight retains a nonprincipal character mode
modulo $3$ on the complete reduced wheel: a bounded product coefficient can
correlate perfectly with the full relaxed survivor count. This refutes the
shortcut that scalar-density centering alone creates Type-II
orthogonality---the exact obstruction Ford--Maynard's framework predicts. The
refutation, the implication from relaxed positivity to prime-plus-$P_2$
production, and the program's exact divisor and bilinear character
properties are proved in
\href{https://github.com/thiagomata/prime-numbers/blob/master/articles/chapter6/relaxed-almost-prime.md}{Relaxed Almost-Prime Production in Sieve Sequences}.

Green and Sawhney's work on prime values of $p^2+nq^2$ demonstrates a modern
successful Type-II strategy using extra algebraic variables, number-field
factorization, and Gowers-norm machinery. Those structural inputs do not
automatically exist for the affine pair $(x,x+2)$. Their result is therefore
a methodological guide, not a theorem that transfers.

The actionable consequence is that any route beyond the exhaustion boundary
of Section~6 must either prove a genuine short-window signed Type-I estimate for
the residue-energy / accepted-boundary quantities of Section~8, or introduce a
new bilinear variable that supplies the Type-II cancellation the affine pair
lacks.
